% GENERATED FILE. Edit canonical lessons, then run npm run generate:books.
% canonical-source-hash: fnv1a64:4d6cf87c228b8262
% canonical-lessons: JA-C194-migi, JA-C194-massugu, JA-C194-magaru, JA-C194-tooru, JA-C194-hakubutsukan

\chapter{Straight On and Turn Right}
\label{ch:ja-straight-on-and-turn-right}
\hlchaptermodality{194}

\begin{chapteropening}
\textbf{By the end of this chapter:} \emph{I can give simple directions.}

Complete the last lesson of chapter 194: I can give simple directions.
\end{chapteropening}

\section[\texorpdfstring{\ja{みぎ} (migi)}{みぎ (migi)}]{\ja{みぎ} (migi) — right; right-hand side}

\label{lesson:JA-C194-migi}



\begin{quote}
Before the new one: say the Japanese for art museum, then the Japanese for to be visible.
\end{quote}

\subsection*{You'll want to know: \ja{みぎ}}
\textbf{\ja{みぎ}} — \emph{migi} — \textquotedblleft{}right; right-hand side\textquotedblright{}.

\begin{quote}
\textbf{\ja{みぎに} \ja{ぎんこうが} \ja{あります。}} — \emph{migi ni ginkō ga arimasu} — \textquotedblleft{}There is a bank on the right.\textquotedblright{}
\end{quote}

The other side is \textbf{\ja{ひだり}}, \textquotedblleft{}left\textquotedblright{}.

\begin{cousinweb}
\textbf{\ja{みぎ}} (\emph{migi}) is a native Japanese word (wago), inherited rather than borrowed from Chinese; its kanji was attached later as a kun'yomi (JMdict, KANJIDIC2).
\end{cousinweb}

\begin{grammarlens}[title={What you've built}]
One more word for everyday life at A2.
\end{grammarlens}

\subsection*{Your turn}
\begin{itemize}
\raggedright
  \item \emph{Say it:} \emph{migi}
  \item \emph{Say it:} \emph{migi}, once more
  \item \emph{Say it:} \emph{migi ni ginkō ga arimasu} — There is a bank on the right.
  \item \emph{Recall:} say the Japanese for art museum, then the Japanese for to be visible, then say \emph{migi} again
\end{itemize}

\begin{tcolorbox}[breakable,colback=teal!4,colframe=teal!35!black,title={Before you move on}]
What does \ja{みぎ} mean? (\textquotedblleft{}Right; right-hand side\textquotedblright{}.) Say it once more.
\end{tcolorbox}
\section[\texorpdfstring{\ja{まっすぐ} (massugu)}{まっすぐ (massugu)}]{\ja{まっすぐ} (massugu) — straight ahead; straight}

\label{lesson:JA-C194-massugu}



\begin{quote}
Before the new one: say the Japanese for to be visible, then the Japanese for right.
\end{quote}

\subsection*{You'll want to know: \ja{まっすぐ}}
\textbf{\ja{まっすぐ}} — \emph{massugu} — \textquotedblleft{}straight ahead; straight\textquotedblright{}.

\begin{quote}
\textbf{\ja{この} \ja{みちを} \ja{まっすぐ} \ja{あるきます。}} — \emph{kono michi o massugu arukimasu} — \textquotedblleft{}I walk straight along this road.\textquotedblright{}
\end{quote}

\begin{cousinweb}
\textbf{\ja{まっすぐ}} (\emph{massugu}) is native Japanese, built from \emph{ma-}, \textquotedblleft{}true, pure\textquotedblright{}, and \emph{sugu}, \textquotedblleft{}straight\textquotedblright{}, both inherited rather than borrowed; JMdict marks it as usually written in kana (JMdict, KANJIDIC2).
\end{cousinweb}

\begin{grammarlens}[title={What you've built}]
One more word for everyday life at A2.
\end{grammarlens}

\subsection*{Your turn}
\begin{itemize}
\raggedright
  \item \emph{Say it:} \emph{massugu}
  \item \emph{Say it:} \emph{massugu}, once more
  \item \emph{Say it:} \emph{kono michi o massugu arukimasu} — I walk straight along this road.
  \item \emph{Recall:} say the Japanese for to be visible, then the Japanese for right, then say \emph{massugu} again
\end{itemize}

\begin{tcolorbox}[breakable,colback=teal!4,colframe=teal!35!black,title={Before you move on}]
What does \ja{まっすぐ} mean? (\textquotedblleft{}Straight ahead; straight\textquotedblright{}.) Say it once more.
\end{tcolorbox}
\section[\texorpdfstring{\ja{まがる} (magaru)}{まがる (magaru)}]{\ja{まがる} (magaru) — to turn; to bend}

\label{lesson:JA-C194-magaru}



\begin{quote}
Before the new one: say the Japanese for right, then the Japanese for straight ahead.
\end{quote}

\subsection*{You'll want to know: \ja{まがる}}
\textbf{\ja{まがる}} — \emph{magaru} — \textquotedblleft{}to turn; to bend\textquotedblright{}.

\begin{quote}
\textbf{\ja{かどを} \ja{みぎに} \ja{まがります。}} — \emph{kado o migi ni magarimasu} — \textquotedblleft{}I turn right at the corner.\textquotedblright{}
\end{quote}

Polite \textbf{\ja{まがります}}, \emph{magarimasu}; a godan verb.

\begin{cousinweb}
\textbf{\ja{まがる}} (\emph{magaru}) is a native Japanese word (wago), inherited rather than borrowed from Chinese; its kanji was attached later as a kun'yomi (JMdict, KANJIDIC2).
\end{cousinweb}

\begin{grammarlens}[title={What you've built}]
One more word for everyday life at A2.
\end{grammarlens}

\subsection*{Your turn}
\begin{itemize}
\raggedright
  \item \emph{Say it:} \emph{magaru}
  \item \emph{Say it:} \emph{magaru}, once more
  \item \emph{Say it:} \emph{kado o migi ni magarimasu} — I turn right at the corner.
  \item \emph{Recall:} say the Japanese for right, then the Japanese for straight ahead, then say \emph{magaru} again
\end{itemize}

\begin{tcolorbox}[breakable,colback=teal!4,colframe=teal!35!black,title={Before you move on}]
What does \ja{まがる} mean? (\textquotedblleft{}To turn; to bend\textquotedblright{}.) Say it once more.
\end{tcolorbox}
\section[\texorpdfstring{\ja{とおる} (tōru)}{とおる (tōru)}]{\ja{とおる} (tōru) — to pass through; to go along, to pass by}

\label{lesson:JA-C194-tooru}



\begin{quote}
Before the new one: say the Japanese for straight ahead, then the Japanese for to turn.
\end{quote}

\subsection*{You'll want to know: \ja{とおる}}
\textbf{\ja{とおる}} — \emph{tōru} — \textquotedblleft{}to pass through; to go along, to pass by\textquotedblright{}.

\begin{quote}
\textbf{\ja{まいにち} \ja{この} \ja{みちを} \ja{とおります。}} — \emph{mainichi kono michi o tōrimasu} — \textquotedblleft{}I go along this road every day.\textquotedblright{}
\end{quote}

Polite \textbf{\ja{とおります}}, \emph{tōrimasu}; a godan verb.

\begin{cousinweb}
\textbf{\ja{とおる}} (\emph{tōru}) is a native Japanese word (wago), inherited rather than borrowed from Chinese; its kanji was attached later as a kun'yomi (JMdict, KANJIDIC2).
\end{cousinweb}

\begin{grammarlens}[title={What you've built}]
One more word for everyday life at A2.
\end{grammarlens}

\subsection*{Your turn}
\begin{itemize}
\raggedright
  \item \emph{Say it:} \emph{tōru}
  \item \emph{Say it:} \emph{tōru}, once more
  \item \emph{Say it:} \emph{mainichi kono michi o tōrimasu} — I go along this road every day.
  \item \emph{Recall:} say the Japanese for straight ahead, then the Japanese for to turn, then say \emph{tōru} again
\end{itemize}

\begin{tcolorbox}[breakable,colback=teal!4,colframe=teal!35!black,title={Before you move on}]
What does \ja{とおる} mean? (\textquotedblleft{}To pass through; to go along, to pass by\textquotedblright{}.) Say it once more.
\end{tcolorbox}
\section[\texorpdfstring{\ja{はくぶつかん} (hakubutsukan)}{はくぶつかん (hakubutsukan)}]{\ja{はくぶつかん} (hakubutsukan) — museum}

\label{lesson:JA-C194-hakubutsukan}



\begin{quote}
Before the new one: say the Japanese for to turn, then the Japanese for to pass through.
\end{quote}

\subsection*{You'll want to know: \ja{はくぶつかん}}
\textbf{\ja{はくぶつかん}} — \emph{hakubutsukan} — \textquotedblleft{}museum\textquotedblright{}.

\begin{quote}
\textbf{\ja{はくぶつかんは} \ja{こうえんの} \ja{なかに} \ja{あります。}} — \emph{hakubutsukan wa kōen no naka ni arimasu} — \textquotedblleft{}The museum is inside the park.\textquotedblright{}
\end{quote}

\begin{cousinweb}
\textbf{\ja{はくぶつかん}} (\emph{hakubutsukan}) is Sino-Japanese (kango): three Chinese characters read with their on'yomi, readings borrowed from Chinese. They are \emph{haku}, which the Guangyun glosses \textquotedblleft{}broad\textquotedblright{} (Middle Chinese \emph{pak}), \emph{butsu}, \textquotedblleft{}thing\textquotedblright{} (\emph{mjut}) and \emph{kan}, \textquotedblleft{}building\textquotedblright{} (\emph{kwanH}) (KANJIDIC2; Guangyun and Middle Chinese in Baxter's transcription, via tshet-uinh).
\end{cousinweb}

\begin{grammarlens}[title={What you've built}]
One more word for everyday life at A2.
\end{grammarlens}

\subsection*{Your turn}
\begin{itemize}
\raggedright
  \item \emph{Say it:} \emph{hakubutsukan}
  \item \emph{Say it:} \emph{hakubutsukan}, once more
  \item \emph{Say it:} \emph{hakubutsukan wa kōen no naka ni arimasu} — The museum is inside the park.
  \item \emph{Recall:} say the Japanese for to turn, then the Japanese for to pass through, then say \emph{hakubutsukan} again
\end{itemize}

\begin{tcolorbox}[breakable,colback=teal!4,colframe=teal!35!black,title={Before you move on}]
What does \ja{はくぶつかん} mean? (\textquotedblleft{}Museum\textquotedblright{}.) Say it once more.
\end{tcolorbox}
